\begin{abcliste}
\abc Im höchsten Punkt muss die Zentripetalbeschleunigung mindestens $g$ betragen:
\begin{align}
 \frac{v^2}{r} &= g \\
 v &= \vAbF = \sqrt{\gE\cdot\rAb} = \vAb \approx \result{\vAbP{0}{2}}
\end{align}

\abc Wir nehmen vereinfachend an, die Bahn fahre den ganzen Kreis mit dieser Geschwindigkeit. Dann dauert eine Runde
\begin{align}
 T &= \frac{2\pi r}{v} = \TAbF = 2\pi\cdot\sqrt{\frac{\rAb}{\gE}} = \TAb \approx \result{\TAbP{0}{2}}
\end{align}
In Wirklichkeit ist die Bahn weiter unten schneller, die Fahrt dauert also etwas weniger lang.
\end{abcliste}
